% ==================================================================
%  Answers to lessons 5 to 9.
% ==================================================================

\subsection*{Lesson 5: what a user sees}

\begin{answer}{absorb}
\begin{enumerate}[leftmargin=*, itemsep=1pt]
\item $\Delta$ enters every code and phase row of system $g$ identically; the
      receiver clock (with the inter-system bias for a second system) takes it up,
      as $\hat a_C$ took up $a_A$ in \cref{q:toy-user}. The geometry, hence the
      position, is unchanged.
\item A constant $d$ enters only that satellite's phase rows, as a constant: the
      float ambiguity moves, $N \to N + d/\lambda$. The position is unchanged; only
      the ambiguity's meaning changes, which matters once a user fixes integers.
\item $(\delta, \mu_2\delta)$ adds $\mu_\sigl\delta$ to code row $\sigl$, the
      ionosphere's own signature, so $I^\sat \to I^\sat + \delta$ absorbs it; on
      phase that shift enters as $-\mu_\sigl\delta$, which the float ambiguities
      absorb. The position is unchanged; the user's slant ionosphere is biased by
      $\delta$. The ionosphere-free combination of $(\delta, \mu_2\delta)$ is zero, so
      an ionosphere-free user sees nothing at all.
\item The per-system mean at each epoch (case 1, the receiver clock), then each
      satellite's mean over the session (case 2, the ambiguity). Case 3 does not
      arise for clocks. The S-system statement: the user uses only the
      block-diagonal part of the correction covariance, ``the rank-one $[e;e]$
      (clock/geometry) and $[-\mu;\mu]$ (iono) parts lie in the user's own design
      range and are absorbed''.
\end{enumerate}
\src{\texttt{SBASIS-THEORY-E.md} lines 158 to 182; \texttt{SBASIS-CONTRACT.md}
lines 98 to 102, 114 to 133; status report, step 5; \texttt{results/t3/results.md}
lines 64 to 66.}
\end{answer}

\begin{answer}{sline}
\begin{center}\footnotesize
\begin{tabular}{@{}p{4.6cm}p{5.4cm}p{4.8cm}@{}}
\toprule
Field & Declares & Defined or written at \\
\midrule
\texttt{S30000SIM} & the run's reference (clock-datum) station; the pins act on satellite clocks & \texttt{ApplicationNetworkSetup.cpp} line 1485 \\
\texttt{1}, \texttt{nominal-wu93} & bias sign convention; phase wind-up model & \texttt{Sbasis.hpp} lines 98 to 108, 215 \\
\texttt{0} (5th field) & ambiguity datum: ``legacy single pivot and prioritized anchor'' & \texttt{Sbasis.hpp} lines 129 to 132 \\
\texttt{0} (6th field) & clock datum: one hard satellite clock pin per constellation & \texttt{Sbasis.hpp} lines 162 to 166 \\
\texttt{1}, \texttt{0} & the network emits its own code biases; code-bias datum id & \texttt{Sbasis.hpp} lines 229, 238 \\
\texttt{0.05}, \texttt{1} & clock discontinuity threshold (m); discontinuities emitted & \texttt{NetworkProductEmitter.hpp} lines 195 to 207 \\
the SP3 and CLK names & the a priori orbit and clock; every C record corrects this clock & \texttt{ApplicationNetworkSetup.cpp} lines 1524 to 1566 \\
\texttt{E:E02}, \texttt{G:G01} & realized clock pivots (C $= 0$ at the first epoch) & \texttt{SolPrep.cpp} lines 1176 to 1179 \\
\texttt{cbjump}, \texttt{atx}, \texttt{oref} & code-bias jump test; antenna file; orbits at the centre of mass & \texttt{NetworkProductEmitter.hpp} lines 222 to 228; \texttt{NetworkProductIO.hpp} lines 270 to 273 \\
\texttt{discfrac}, \texttt{hsig} & phase jump gate of 0.25 cycle; $\sigma_0$ of the H records & \texttt{NetworkProductEmitter.hpp} lines 238 to 248; \texttt{NetworkProductIO.hpp} lines 238 to 247 \\
\texttt{cbpin=...} & the ionosphere-free pin pair per system & \texttt{SolPrep.cpp} lines 1118 to 1157 \\
\texttt{sat\_gf\_datum=broadcast\_tgd} & the geometry-free datum at $c\,T_{GD}$ & \texttt{SolPrep.cpp} lines 1165 to 1172 \\
\texttt{pbhold=2} & the phase-clock hold (\cref{q:hold}) & \texttt{ApplicationNetworkSetup.cpp} lines 1599 to 1601 \\
\bottomrule
\end{tabular}
\end{center}
The reference choices of \cref{sec:l4} are the two datum ids, the pivots, the
\texttt{cbpin} pairs and \texttt{sat\_gf\_datum}; \texttt{pbhold} declares how the
phase biases are held. A current run appends \texttt{wumode=accumulated}. The
real-data S line declares only \texttt{cbpin=G:C1C-C2W}: the Galileo pairs are
``NOT declared (per-satellite split)'' (\texttt{runs/\_rf0hv15\_field6/stdout.txt}
line 686).
\src{\texttt{src/core/products/NetworkProductIO.cpp} lines 465 to 602 and the lines
in the table (paths under \texttt{src/core/network/} and \texttt{include/}).}
\chk
\begin{cmd}
sed -n '2p' D:/GnssProcessorSupport/runs_nom/_rap_scb_g_q6/products.txt
\end{cmd}
\end{answer}

\begin{answer}{step}
\begin{enumerate}[leftmargin=*, itemsep=1pt]
\item The awk prints line 134213, \texttt{E,2313,271230.00000715256,400,0,0,1,E}
      (Galileo, 03:20:30), and line 189631, \texttt{E,2313,276149.99999547005,564,0,0,1,G}
      (GPS, 04:42:30). The script prints ``31 satellites, mean step 1.6278 m, sd
      after the mean 0.0027 m''; an ordinary pair of epochs (276090 to 276120) gives
      $-0.0001$ and 0.0008~m.
\item The handover of \cref{q:toy-setchange} (3): the pivot moved from G01 to G02 and
      every GPS clock stepped by the old pivot-relative value. A user's receiver
      clock absorbs the common 1.63~m; 3~mm per satellite stays visible.
\item A raw-step grade sees 31 satellites jump 1.6~m at once; the datum-invariant
      step is 3~mm. On real data the GPS group mean steps by about $+154$~cm at its
      pivot change. \emph{Checked.}
\end{enumerate}
\src{\texttt{runs/\_rf0hv15\_sim6/products.txt}; \texttt{results/t3/results.md}
lines 105 to 107.}
\end{answer}

\begin{answer}{refsat}
\begin{enumerate}[leftmargin=*, itemsep=1pt]
\item A chord is $e_\sat = s_\sat - s_{\mathrm{ref}}$: $\mathrm{Var}(e_\sat) =
      \sigma^2 + \sigma_{\mathrm{ref}}^2$ and $\mathrm{Cov}(e_\sat, e_k) =
      \sigma_{\mathrm{ref}}^2$, so two chords correlate by
      $\sigma_{\mathrm{ref}}^2/(\sigma^2 + \sigma_{\mathrm{ref}}^2)$, 0.5 when
      $\sigma_{\mathrm{ref}} = \sigma$.
\item L1C over all epochs: as traced, median $|e|$ 0.304 cycle, 65.3\,\% round to
      the right integer; reference-free 0.122 cycle, 81.6\,\%. Last quarter: 0.260
      cycle and 69.5\,\% against 0.092 cycle and 86.5\,\%.
\item The reference satellite's error and its returns enter every chord, so a grade
      on the traced chords reports that one element's share: ``Quote DD grades
      ref-free from now on.'' A second trap of the same card: subtracting a group's
      own circular median shrinks uniform noise from 0.251 to 0.145 cycle, so a
      wrapped grade is quoted as a gap to a size-matched null. \emph{Checked.}
\end{enumerate}
\src{\texttt{docs/cards/sim-bed/DD-TRACE-REF-SAT.md} lines 41 to 48, 107 to 110, 143
to 159; \texttt{CLAUDE.md}, ``Grading traps that still apply''.}
\end{answer}

\subsection*{Lesson 6: the satellite clock}

\begin{answer}{qstep}
\begin{enumerate}[leftmargin=*, itemsep=1pt]
\item $\sqrt{2.9\times10^{-8}\times30} = 0.933$~mm per epoch;
      $\sqrt{2.9\times10^{-8}\times21{,}600} = 2.50$~cm over 6~h. The value is
      measured: ``$q_{\mathrm{phys}} \approx 2.9\times10^{-8}$~m$^2$/s $\to$ 0.94~mm
      per 30~s''.
\item The drift after $n$ epochs is $0.30n$~mm, the walk's $1\sigma$ $0.933\sqrt n$~mm;
      they are equal at $n = (0.933/0.30)^2 = 9.7$ epochs, about 5~min. At 6~h the
      drift is 21.6~cm, 8.6 times the walk. A step standard deviation is being
      compared with a signed, persistent increment: a drift grows as $n$, a random
      walk's spread as $\sqrt n$, so a walk of this density cannot follow a sustained
      rate. Card SCB reads the outcome as partial tracking.
\end{enumerate}
\src{\texttt{config/golden/golden\_sim\_refound.json} line 323;
\texttt{include/network/solve/ClockAnchor.hpp} lines 35 to 36;
\texttt{docs/cards/clock/Q3-PHYSICAL-FLOOR.md} lines 29 to 32; card SCB lines 29 to 35.}
\end{answer}

\begin{answer}{tau}
\begin{enumerate}[leftmargin=*, itemsep=1pt]
\item $R = 0.0923^2/2 = 4.26\times10^{-3}$~m$^2$; $K = \sqrt{8.7\times10^{-7}/4.26\times
      10^{-3}} = 0.0143$; $\Delta t/K \approx 2{,}100$~s, about 35~min (the exact
      steady state of the scalar filter gives 2,084~s).
\item $10^{-5}\times2{,}100 \approx 2.1$~cm. Card SCB measured 2,127~s (2,000~s GPS
      only): the anchor's $q$ sets the time constant, within about 2\,\%. Two
      cautions: the argument is scalar and stationary, and the real filter is
      neither; and the grader's error is estimate minus truth with a positive
      coefficient, which a pure lag would make negative. The card calls it a lag and
      does not explain the sign, so the safe wording is ``an error proportional to
      the drift rate, with a coefficient of about 35~min''.
\item $K = \sqrt{1.5\times10^{-4}/4.26\times10^{-3}} = 0.19$, about 160~s (exact
      steady state 146~s), the scale of the measured 52~s: the regime's larger $q$
      follows the drift within the epoch noise. \emph{Checked (arithmetic).}
\end{enumerate}
\src{\texttt{docs/cards/clock/Q3-ANCHOR-DESIGN.md} lines 47 to 56; card SCB lines 124,
128; \texttt{docs/cards/raim-qc/DAY129-FRAME-ANATOMY.md} lines 139 to 143;
\texttt{scripts/simnet\_clock\_truth\_grade.py} line 151.}
\chk
\begin{cmd}
C:/Python313/python.exe - <<'EOF'
import math
R = 0.0923**2 / 2
for q in (2.9e-8, 5e-6):
    K = math.sqrt(q * 30 / R)
    print(q, 'K %.4f' % K, 'tau %.0f s' % (30 / K), 'error %.4f m' % (1e-5 * 30 / K))
EOF
\end{cmd}
\end{answer}

\begin{answer}{hold}
\begin{enumerate}[leftmargin=*, itemsep=1pt]
\item $C - P = c(t) - p(t_0) - [\tilde c(t) - \tilde c(t_0)] = (c - p)(t_0) + M(t) -
      M(t_0)$: constant per satellite up to the system mean $M$, which the receiver
      clock absorbs. On code the user applies $a(t) + C(t)$, the network's estimate;
      on phase it sees $a(t) + C(t) - P(t) = a(t)$ plus a constant: the a priori
      clock's motion. With truth $a(t) + e(t)$, the phase rows are wrong by
      $-[e(t) - e(t_0)]$ plus a constant; the float ambiguity absorbs the constant,
      the change passes one to one.
\item Mean square over hour 6 ($t$ from 18,000 to 21,600~s): rate part
      $\sigma_r^2(t_{\mathrm{hi}}^3 - t_{\mathrm{lo}}^3)/(3\times3600) = 10^{-10}\times
      3.931\times10^8 = 0.0393$~m$^2$; walk part $q(t_{\mathrm{lo}} + t_{\mathrm{hi}})/2
      = 0.0006$~m$^2$; rms 0.200~m, against the measured 0.201~m (GPS) and 0.206~m
      (Galileo).
\item The loader refuses \texttt{network.emit\_pb\_hold\_phase\_clock} with
      \texttt{network.apriori\_source} \texttt{sp3\_brdc\_clock}: ``the phase-clock hold
      is valid only with a precise a priori clock product''; with \texttt{brdc\_rts}
      it only warns (lines 2612 to 2617). Part 1 is the reason: the hold hands the
      user the a priori clock's motion, so the a priori clock must be precise. The S
      line declares \texttt{pbhold=2} (3 for the smoothed hold of the self-contained
      postures). On real data, ignoring the precondition cost the user 1.191 and
      0.743~m in hours 1 and 2, against 0.108 and 0.024~m with CODE.
      \emph{Checked.}
\end{enumerate}
\src{\texttt{include/network/emit/PbHold.hpp} lines 60 to 66; card LAH lines 30 to 36;
card SCB lines 129 to 145; \texttt{src/core/config/GnssConfigLoadNetwork.cpp} lines 311
to 319, 2612 to 2617; card SCL lines 22 to 26, 76 to 78.}
\chk
\begin{cmd}
C:/Python313/python.exe - <<'EOF'
lo, hi = 18000.0, 21600.0                      # hour 6, s after the reset
rate = 1e-10 * (hi**3 - lo**3) / (3 * 3600)     # rate part, m^2
walk = 3e-8 * (lo + hi) / 2                     # random-walk part, m^2
print((rate + walk) ** 0.5)
EOF
sed -n '311,319p' src/core/config/GnssConfigLoadNetwork.cpp
\end{cmd}
\end{answer}

\begin{answer}{postures}
Reference case \texttt{\_rap\_scb\_g\_k6}: clock level at hour 6 0.0533~m;
per-satellite mean over hours 2 to 6, rms 0.0426~m over 55 satellites; correlation
with the drawn rate $+0.524$; user by hour 0.1976, 0.2012, 0.3189, 0.2142, 0.2370,
0.3376~m. Self-contained regime \texttt{\_rap\_scb\_c\_k6}: 0.0232~m; 0.0205~m over
31 satellites; $+0.026$; user 0.2287, 0.1296, 0.1003, 0.0211, 0.0311, 0.0260~m. The
regime follows the drift (no correlation with the rate); the reference case does
not. Its user error is about ten times its clock error because the hold passes the
a priori clock's error to the phase rows whatever the estimator tracks
(\cref{q:hold}). Not shown: one realization of an error with chosen statistics; the
sign question of \cref{q:tau}; the regime case is GPS only (31 satellites against
55); neither the regime nor the native posture is armed. Card SCB's text says the
regime and native users ``do not move'' between the two data sets; its tables show
the native user up to 2.1~cm higher on the clock data set in hours 1 to 3.
\emph{Checked.}
\src{the two grader outputs; card SCB lines 8, 10, 122 to 145.}
\end{answer}

\begin{answer}{nbp}
\begin{enumerate}[leftmargin=*, itemsep=1pt]
\item $0.0024 - 0.0007 + 0.0099 = 0.0116$~m, not 0.0171~m: the steps add exactly for
      each station-day, but the median of a sum is not the sum of the medians. The
      clock's information is the broadcast-clock step, about 1.0~cm.
\item The native posture against the GPS-only reference configuration of
      2026-10-07 (GPS and Galileo since 2026-10-13), days 130 to 132 (9 to 11 May
      2024), 16 stations, as the median paired change of each station's hours 2 to 6
      error; 19 of 48 station-days (39.6\,\%) are worse by more than 2~cm; day 129 is
      its most favourable day ($+0.3$~cm); the first hour adds $+7.65$~cm.
\end{enumerate}
\src{card NBP lines 170 to 177, 184 to 190, 210 to 219; card GAD.}
\end{answer}

\subsection*{Lesson 7: ambiguities}

\begin{answer}{nl}
$\lambda_1 = c/f_1 = 0.190294$~m, $\lambda_2 = 0.244210$~m, $\mu_2 = (77/60)^2 =
1.646944$. $\cos\theta = (\lambda_1 + \mu_2\lambda_2)/(|(1,\mu_2)|\,|(\lambda_1,
\lambda_2)|) = 0.592494/(1.926765\times0.309597) = 0.99325$, $\theta = 6.66^\circ$.
The part orthogonal to the ionosphere is $0.3096\times\sin6.66^\circ = 0.0359$~m of
the 0.31~m step. ``99.3\,\%'' is the cosine, a share of amplitude; as a share of
variance ($\cos^2$) it is 98.7\,\%. The metric assumes equal phase noise in metres on
both bands. \emph{Checked.}
\src{card ARL line 43 (the statement); the arithmetic.}
\chk
\begin{cmd}
C:/Python313/python.exe - <<'EOF'
import math
c = 299792458.0
l1, l2 = c / 1575.42e6, c / 1227.60e6
mu2 = (1575.42 / 1227.60) ** 2
k = (l1 + mu2 * l2) / math.hypot(1, mu2) / math.hypot(l1, l2)
print(k, math.degrees(math.acos(k)), math.hypot(l1, l2) * math.sin(math.acos(k)), k * k)
EOF
\end{cmd}
\end{answer}

\begin{answer}{cycles}
\begin{enumerate}[leftmargin=*, itemsep=1pt]
\item Cycles per metre of L1 ionosphere: L1 $1/\lambda_1 = 5.255$, L2 $\mu_2/\lambda_2
      = 6.744$, wide lane $5.255 - 6.744 = -1.489$ (the card: ``L1C 5.255, L2W
      6.744''). $d = 0.4233$~m: L1 2.22 cycles, wide lane 0.63 cycle. $d = 0.0419$~m:
      L1 0.22 cycle, wide lane 0.062 cycle.
\item 0.1 L1 cycle $= 0.0190$~m; 1~TECU $= 40.3\times10^{16}/f_1^2 = 0.1624$~m of L1
      delay, so 0.1 cycle is 0.117~TECU. The no-pin error, 0.4233~m $= 2.61$~TECU, is
      22 times that: ``2.61 TECU, 22$\times$ the 0.117 TECU that moves a fraction''.
      The 0.4233~m is a per-line-of-sight rms, so a double difference of several such
      terms can be larger. \emph{Checked.}
\end{enumerate}
\src{\texttt{docs/cards/raim-qc/T3-IONO-MECHANISM.md} lines 326 to 331;
\texttt{docs/SPEC\_oracle\_pin.md} lines 514 to 517, 541 to 546.}
\end{answer}

\begin{answer}{pins}
\begin{enumerate}[leftmargin=*, itemsep=1pt]
\item The pins make the network emit $\bar b^\sat = b^\sat + N^\sat_A - \zeta^\sat$
      (in cycles), and the user then estimates $\tilde N^\sat_C = N^\sat_C - N^\sat_A
      + \zeta^\sat = (z^\sat_C - z^\sat_A) + (\beta_C - \beta_A) + \zeta^\sat$.
\item $\tilde N^\sat_C - \tilde N^k_C = [(z^\sat_C - z^k_C) - (z^\sat_A - z^k_A)] +
      (\zeta^\sat - \zeta^k)$: the bracket is an integer double difference (the
      receiver phase biases cancel), so the difference is an integer exactly when
      $\zeta^\sat - \zeta^k$ is. Integer pins (\cref{q:toy-user}: $\zeta = 0$) give
      integers; float pins give a fraction.
\item ``The WIRE's gauge is an arbitrary receiver's FLOAT ambiguity'': the anchors
      are float targets, re-pinned at the successor's float on every handover, so the
      user's between-satellite differences carry fractions spread over a whole cycle,
      and the published biases sit at the uniform null (\cref{q:null}). Card IRD: a
      user's single difference is integer if and only if every emitted bias
      difference is; it proposes to pin a spanning tree of each observation graph
      ($n_c + m_c - 1$ edges) to integers plus one real-valued root receiver
      phase-bias pin, so that every chord is an integer (after de
      Jonge~\cite{deJonge1998processing}). Not built (IRD-1 is the next step).
\end{enumerate}
\src{\texttt{docs/network\_ambiguity\_resolution\_design.md} lines 335, 352 to 356; card
IRD lines 19, 41 to 49.}
\end{answer}

\begin{answer}{oracle}
Row 1b, the double-difference-invariant per-line-of-sight L1 slant ionosphere
error, carries the deficit: 0.4233~m with no pin, 0.0419~m with every ambiguity
pinned (10.1 times smaller; 2.61 to 0.26~TECU). ``The \texttt{none} column's row 1b
IS the integer-information deficit''. The raw per-line-of-sight error (2.73~m) is
about 97\,\% station and satellite lumps, which the reference choice and the user
absorb, so it cannot measure integer information. Partial pins do not interpolate:
L5Q alone 0.2793~m (better); L1C with L2W 0.5632~m (worse than no pin, although every
double difference is then integer); one station 0.4476~m; one satellite 0.4741~m.
``Integer information is worth having only when it is complete on the LOS.'' The
proposed cause (the bands of a line of sight share one ionosphere state, so exact
integers on a subset over-determine it) is called a consistent reading, not a proof.
The numbers come from a 60-epoch warm resume on \texttt{simnet60\_quiet\_wu93\_24h},
a data set whose numbers \texttt{CLAUDE.md} marks superseded: the reading stands, the
values are not comparable with today's simulated data.
\src{\texttt{docs/SPEC\_oracle\_pin.md} lines 20 to 29, 193 to 195, 514 to 517, 534 to
546, 558 to 568, 611 to 627, 654 to 663.}
\end{answer}

\begin{answer}{null}
First command (4~s): the published GPS phase biases' wide-lane median $|$fraction$|$
by hour from 00:30:15 is 0.254, 0.260, 0.259, 0.255, 0.253, 0.252, the narrow lane
0.264 to 0.266, against a uniform null of 0.250 (band about 0.22 to 0.27). Second (4~s):
``CORRECT-FIX RATE among kept: 5572/5572 = 100.00\%'', ``WRONG FIXES: none''. The
integers are in the network's float solution (a search on its float states and
covariance fixes 5,572 of 13,138 settled double differences of that case, all
correct), but the published phase biases carry none of it: they sit at the uniform
null, as \cref{q:pins} predicts for float pins. A user cannot fix integers from these
corrections. The 42\,\% case is a simpler 2~h simulation without time-correlated phase
errors or slips; with the time-correlated phase error the share is 18.5\,\%.
\emph{Checked.}
\src{the two outputs; card ARL lines 98 to 102, 142.}
\end{answer}

\subsection*{Lesson 8: simulated data}

\begin{answer}{quiet}
The quiet data set's truth clock is the CODE final clock, and the network takes the
same product as its a priori clock, so the estimable satellite clock truth is
identically zero (card SCB, section 1). \texttt{results/t2/results.md} confirms it:
no \texttt{sat\_clock} key in the generator's configurations, no ARMED line, no
\texttt{truth\_satclock.csv}. The two \texttt{ls} commands print 0 and 61. The
0.9~cm agreement therefore ``certifies the algebra and the biases, not the clock'':
a simulation grades what it injects, plus contrasts in which every reused term
cancels. \emph{Checked.}
\src{card SCB lines 17 to 25; \texttt{results/t2/results.md} lines 123 to 126;
\texttt{docs/cards/sim-bed/SIMNET-REF.md} lines 269 to 272.}
\end{answer}

\begin{answer}{start}
\begin{enumerate}[leftmargin=*, itemsep=1pt]
\item 0.504167~h $= 1{,}815.0012$~s; epoch $k$ is $30k$~s after the first, and the
      first $k$ with $30k \geq 1{,}815.001$ is 61 (1,830~s): tow 261060, 00:31:00.
      0.4958333333~h $= 1{,}785$~s gives $k = 60$: tow 261030, 00:30:30 (the T2
      restart). Each offset sits 15~s, half an epoch, below the epoch it selects, so
      rounding cannot move the start by an epoch. The grading origin 00:30:15 (tow
      261015) is not a start epoch: the reference case starts at 00:00:30 and is
      graded from 00:30:15; a case with the offset starts at 00:31:00.
\item Output: 259230 367, 259260 366, 259290 366, 259320 367, 259350 309 (warnings per
      epoch): 5 epochs logged. A warning needs more than 150~s outside the window,
      which opens at tow 259500, so 259380 to 259470 are extrapolated without one: 9
      epochs, 00:00:30 to 00:04:30. ``Faulty in the first 30 minutes'' described the
      grading origin, not the fault; what lasts is the state those epochs leave
      (\cref{q:basis}). The report now says 9 epochs. \emph{Checked.}
\end{enumerate}
\src{\texttt{include/math/TimeUtils.hpp} lines 416 to 448;
\texttt{docs/cards/queues/QPR-ITEMS-5-8.md} lines 239, 256;
\texttt{src/core/model/ModelGeometry.cpp} line 1204;
\texttt{docs/refound/CTU2\_level\_recursion\_followup\_2026-09-29.md} lines 29 to 43.}
\chk
\begin{cmd}
head -2 D:/GnssProcessorSupport/runs_nom/_rap_scb_g_q6/logs/closed_loop_trace.csv
\end{cmd}
\end{answer}

\begin{answer}{basis}
Baseline (from 00:00:30): level by hour 0.7087, 0.3487, 0.3303, 0.3157, 0.3096,
0.3108~m; per-satellite mean over hours 2 to 6, rms 0.3179~m over 55 satellites
(G06 $-1.198$, E05 $+0.872$, E36 $-0.846$). Card SCB's quiet case (from 00:31:00):
level 0.0809, 0.0291, 0.0226, 0.0202, 0.0193, 0.0185~m; rms 0.0186~m. With
\texttt{--no-ds} the baseline's rms reads 0.7784~m. The datum share is the
ionosphere-free combination of the truth code biases over the pinned pair: the
network's clock contains it by construction of the code-bias datum
(\cref{q:cbdatum}), so it belongs to the truth expressed in the network's basis.
Truth is transformed; differencing the estimates instead would bias the variance
factor downward, because the estimates share the datum. The baseline's user does not
show the 0.32~m: a float user cannot see a per-satellite constant (\cref{q:absorb}).
\emph{Checked.}
\src{the grader outputs; card SCB lines 87, 102, 107 to 109;
\texttt{scripts/simnet\_clock\_truth\_grade.py} lines 138 to 153; \texttt{SIMNET-REF.md}
lines 276 to 283.}
\end{answer}

\begin{answer}{floor}
Ours minus perfect, in metres: $+0.0030$, $+0.0026$, $+0.0085$, $-0.0007$, $+0.0005$,
$+0.0002$. The largest cost is 0.85~cm, in hour 3 (the report's ``at most 0.9~cm'');
hours 4 to 6 show nothing measurable. The floor (perfect corrections, hours 4 to 6)
is 1.73, 1.81 and 1.76~cm, in the height: median $|U|$ 1.54, 1.54, 1.65~cm against
$|E| \leq 0.46$~cm and $|N| \leq 0.55$~cm. Hour 6 is 30~min long, because the network
run ends at 06:00.
\src{\texttt{results/t2/results.md} lines 195 to 206; \texttt{results/t2/metrics.csv}.}
\end{answer}

\subsection*{Lesson 9: real data}

\begin{answer}{census}
\texttt{rapid\_arm.py score}: \texttt{\_rf0hv15\_field6} median by hour 0.06, 0.01,
0.02, 0.01, 0.02, 0.01~m, 688 epochs below 10~cm; \texttt{\_rf0hv15\_sim6} 0.14, 0.07,
0.03, 0.02, 0.03, 0.02~m, 595. \texttt{gcensus.py}: \texttt{galp} (GPS and Galileo
corrections), 16 stations, median by hour 0.136, 0.054, 0.041, 0.036, 0.040, 0.037~m,
WTZZ 0.058, 0.011, 0.018, 0.010, 0.015, 0.011~m; paired against \texttt{galpc} (GPS
only), hours 2 to 6 median $-0.0198$~m, 10 better, 2 worse (HARB $+0.0158$, DAV1
$+0.0605$), hour 1 $-0.2203$~m; \texttt{galpg} (GPS-only users of the GPS and Galileo
corrections) $-0.0000$~m, ARHT $+0.1958$~m. Better and worse mean a change of more
than 1~cm. The first hour: WTZZ00DEU 5.8 against a mean of 1.3~cm over hours 2 to 6
(4.5 times); the median station 13.6 against 4.2~cm (3.3 times). Mode:
\texttt{position\_q} 100 (moving receiver), \texttt{ar\_mode} 0 (no integer fixing).
None of the 17 user stations is in the 58-station network; WTZZ00DEU shares its site
with network stations WTZR00DEU and WTZS00DEU (``the colocated consumer hold-out'';
its RINEX header position is about 3~m from WTZR00DEU's). Whether that explains its
rank has not been tested. \emph{Checked.}
\src{the outputs; \texttt{runs/\_rf0hv15\_field6/arm.json}; \texttt{scripts/clock\_level/gcensus.py}
lines 151 to 153; card GAD lines 70 to 75; \texttt{docs/refound/RF-PD\_station\_rejection\_2026-09-15.md}
line 87.}
\chk
\begin{cmd}
C:/Python313/python.exe - <<'EOF'
import json
cfg = json.load(open('config/golden/golden_field_refound.json'))
names = [s['name'] if isinstance(s, dict) else s for s in cfg['network']['stations']]
print(len(names), [n for n in names if n.startswith('WTZ')])
EOF
\end{cmd}
\end{answer}

\begin{answer}{fivedays}
Hours 2 to 6 of days 128 to 132: the 25 sixteen-station medians have a median of
0.038~m and a mean of 0.0406~m; the 25 WTZZ00DEU values a median of 0.018~m. On day
130 WTZZ00DEU reads 3.9, 5.9 and 5.3~cm in hours 4 to 6 against the median station's
3.5, 3.2 and 3.8~cm: ``the best station'' holds on 8 May 2024 (day 129), not on every
day. The report, which first quoted 3.9~cm for the 16 stations, now quotes 3.8~cm
(mean 4.1~cm). \emph{Checked.}
\src{\texttt{D:/GnssProcessorSupport/runs\_nom/mday\_galileo\_grades.txt}.}
\chk
\begin{cmd}
C:/Python313/python.exe - <<'EOF'
import re, statistics as st
cen, wtz = [], []
for l in open('D:/GnssProcessorSupport/runs_nom/mday_galileo_grades.txt'):
    m = re.match(r'== mdg\d+c: .*by hour ([\d. ]+); WTZZ ([\d. ]+)', l)
    if m:
        cen += [float(x) for x in m.group(1).split()[1:]]
        wtz += [float(x) for x in m.group(2).split()[1:]]
print(len(cen), st.median(cen), round(st.mean(cen), 4), st.median(wtz))
EOF
\end{cmd}
\end{answer}

\begin{answer}{t3}
Against the GFZ final clock, after the per-system epoch mean and each satellite's
mean are removed (STD): our network (CODE final plus c0) GPS 4.3 and Galileo 7.6~cm
over 01:00 to 06:00 (6.4 and 11.6~cm including the first hour; Galileo RMS 15.8 and
25.3~cm because of E12); CODE final 1.4 and 2.6~cm over 00:00 to 06:00, 1.2 and 2.2~cm
over 01:00 to 06:00, 1.9 and 3.3~cm over 13:24 to 24:00; the IGS real-time service 4.2
and 3.2~cm over 13:24 to 24:00, the only hours archived. GFZ is the reference because
the network's a priori clock is CODE's: against CODE the network would be measured
partly against itself. Over the same hours the corrected clock is about 3.5 times
noisier than CODE alone for GPS. The correlation of c0 with (GFZ minus CODE) is GPS
$-0.08$, Galileo $+0.01$, and with (ESA minus CODE) GPS $-0.18$, Galileo $+0.01$: c0
carries no information about CODE's error; relative to its a priori it behaves as
added noise of 6.3~cm (GPS). A clock comparison alone does not isolate what a user
obtains: the user's phase term is phase bias minus clock, which the hold keeps at
CODE's clock (\cref{q:hold}). \emph{Checked.}
\src{\texttt{results/t3/table\_clock\_report.tex}; \texttt{results/t3/results.md}
lines 9 to 14, 100 to 113; \texttt{results/t3/clock\_stats.csv}, rows
\texttt{COD,GFZ,01-06}.}
\chk
\begin{cmd}
T3=D:/GnssProcessorSupport/presentation/status_update_2026-10-29/results/t3
grep "COD,GFZ,01-06" $T3/clock_stats.csv
\end{cmd}
\end{answer}

\begin{answer}{capstone}
\begin{enumerate}[label=(\alph*), leftmargin=*, itemsep=2pt]
\item \emph{Model and check.} May: user and network use one undifferenced,
      uncombined observation model, and the user one filter in every step; the rank
      gate found exactly the declared 2 directions (one clock per constellation) on
      720 of 720 epochs, none withheld, on both reference cases (\cref{q:rankgate}).
      May not: that the network uses the user's filter (it has its own
      information-form solve); that the check proves the choices harmless (it counts
      null directions only); that it covers each station's level (\cref{q:level}).
\item \emph{Reference choices.} May: they are declared in the correction file and
      every change of pivot is marked (\cref{q:sline,q:step}). May not: ``a few''
      (226 rows and 120 prior entries at the first simulated epoch,
      \cref{q:rows}); all declared on real data (the Galileo pin pairs are not); that
      they keep a user's integers (\cref{q:pins}).
\item \emph{Satellite clock.} May: its precision comes from CODE's final clock; the
      network's correction adds code noise and no clock information
      (\cref{q:t3}); on simulated clock error the self-contained postures reach
      about 2~cm, unarmed (\cref{q:postures}). May not: that it beats real-time
      clocks; that the network is free of outside clock products; that the quiet
      simulation tests the clock (\cref{q:quiet}). Without CODE's clock a user was
      1.7~cm worse at the median, against an older GPS-only configuration
      (\cref{q:nbp}).
\item \emph{Integers.} May: on a simplified simulation a log-only search fixed 42\,\%
      of 13,138 double differences, none wrongly (\cref{q:null}). May not: integers
      in the corrections; integers on real data; producer-side fixing, whose standing
      verdict is DO NOT BUILD (\texttt{docs/network\_ambiguity\_resolution\_design.md}
      lines 3 to 5).
\item \emph{Results.} May: on 8 May 2024, WTZZ00DEU 1.0 to 1.8~cm and the median of
      16 further stations 3.6 to 5.4~cm in hours 2 to 6, without integers; over five
      days the median station 3.8~cm (mean 4.1~cm); on the quiet simulation within
      0.9~cm of perfect corrections (\cref{q:census,q:fivedays,q:floor}). May not:
      the best station on every day; WTZZ00DEU as a user unrelated to the network
      (it shares its site with two network stations); claims from the first hour.
\end{enumerate}
These boundaries were applied to the status report's wording on 24~October 2026.
\src{the answers named.}
\end{answer}
